\documentclass[12pt]{article}
\usepackage{amsmath,amssymb,amsfonts}
\usepackage[margin=1in]{geometry}

\begin{document}

\textbf{Chapter 6, Problem 8.} Using the results of Problem 6.6, show that $\log z$, $z^{1/n}$, and $z^{1/3}$ are analytic everywhere on their appropriate Riemann sheets except at 0.

\bigskip
\textbf{Solution.}

From Problem 6 (the polar form of the Cauchy--Riemann equations):
\[
\frac{\partial u}{\partial r} = \frac{1}{r}\frac{\partial v}{\partial \theta}, \qquad \frac{\partial v}{\partial r} = -\frac{1}{r}\frac{\partial u}{\partial \theta}.
\]

\textbf{(i) $\log z$:} On a given Riemann sheet, $\log z = \ln r + i\theta$ where $\theta$ ranges over an interval of length $2\pi$. So $u = \ln r$ and $v = \theta$. Check:
\[
\frac{\partial u}{\partial r} = \frac{1}{r}, \qquad \frac{1}{r}\frac{\partial v}{\partial \theta} = \frac{1}{r}. \quad \checkmark
\]
\[
\frac{\partial v}{\partial r} = 0, \qquad -\frac{1}{r}\frac{\partial u}{\partial \theta} = 0. \quad \checkmark
\]

The derivative is $f'(z) = u_r + iv_r$ times $e^{-i\theta}/1$... more precisely, $f'(z) = e^{-i\theta}(u_r + iv_r) = e^{-i\theta}\cdot(1/r + i\cdot 0) = 1/(re^{i\theta}) = 1/z$.

This is defined for all $r > 0$, i.e., for all $z \neq 0$.

\textbf{(ii) $z^{1/n}$:} On sheet $k$, $z^{1/n} = r^{1/n}e^{i(\theta+2k\pi)/n}$. So:
\[
u = r^{1/n}\cos\!\frac{\theta+2k\pi}{n}, \qquad v = r^{1/n}\sin\!\frac{\theta+2k\pi}{n}.
\]

Let $\phi = (\theta + 2k\pi)/n$. Then:
\[
\frac{\partial u}{\partial r} = \frac{1}{n}r^{1/n-1}\cos\phi, \qquad \frac{1}{r}\frac{\partial v}{\partial \theta} = \frac{1}{r}\cdot r^{1/n}\cdot\frac{1}{n}\cos\phi = \frac{1}{n}r^{1/n-1}\cos\phi. \quad\checkmark
\]
\[
\frac{\partial v}{\partial r} = \frac{1}{n}r^{1/n-1}\sin\phi, \qquad -\frac{1}{r}\frac{\partial u}{\partial \theta} = -\frac{1}{r}\cdot r^{1/n}\cdot\left(-\frac{1}{n}\sin\phi\right) = \frac{1}{n}r^{1/n-1}\sin\phi. \quad\checkmark
\]

The Cauchy--Riemann equations are satisfied for all $r > 0$. The derivative is:
\[
f'(z) = e^{-i\theta}(u_r + iv_r) = e^{-i\theta}\frac{1}{n}r^{1/n-1}e^{i\phi} = \frac{1}{n}r^{1/n-1}e^{i(\phi-\theta)}.
\]
Since $\phi - \theta = (\theta+2k\pi)/n - \theta = \theta(1/n-1) + 2k\pi/n$, this simplifies to $\frac{1}{n}z^{1/n-1}$.

This exists and is finite for all $r > 0$, i.e., $z \neq 0$. At $z = 0$ the derivative does not exist (for $n > 1$, the function is not differentiable at the branch point).

The case $z^{1/3}$ follows identically with $n = 3$.

Therefore, $\log z$, $z^{1/n}$, and $z^{1/3}$ are analytic on their respective Riemann sheets for all $z \neq 0$.

\end{document}
